% =====================================================================
%  Tri grossier du corpus : quels critères globaux sont robustes,
%  et avec quels modèles locaux (media_restorer)
%
%  Compilation :  pdflatex rapport-tri-grossier-corpus.tex   (× 3)
% =====================================================================
\documentclass[11pt,a4paper]{article}

\usepackage[utf8]{inputenc}
\usepackage[T1]{fontenc}
\usepackage{lmodern}
\usepackage[french]{babel}
\usepackage[margin=2.4cm]{geometry}
\usepackage{microtype}
\usepackage{parskip}
\usepackage{booktabs}
\usepackage{tabularx}
\usepackage{enumitem}
\usepackage{xcolor}
\usepackage{fancyhdr}
\usepackage{titlesec}
\usepackage[hidelinks]{hyperref}
\usepackage{xurl}             % césure des URL n'importe où

\definecolor{accent}{HTML}{1F4E79}
\definecolor{encadre}{HTML}{F2F5F8}
\definecolor{alerte}{HTML}{A33A2E}
\definecolor{vert}{HTML}{2E6B3E}
\hypersetup{
  colorlinks=true, linkcolor=accent, urlcolor=accent, citecolor=accent,
  pdftitle={Tri grossier du corpus — critères robustes et modèles locaux},
  pdfsubject={media\_restorer},
  pdfkeywords={tri, classification, dessins, WD14, CLIP, DigiKam},
}

\titleformat{\section}{\Large\bfseries\color{accent}}{\thesection}{0.7em}{}
\titleformat{\subsection}{\large\bfseries}{\thesubsection}{0.6em}{}

\pagestyle{fancy}
\fancyhf{}
\fancyhead[L]{\small\color{accent}media\_restorer — tri grossier}
\fancyhead[R]{\small\thepage}
\renewcommand{\headrulewidth}{0.4pt}

\newcommand{\aretenir}[1]{%
  \par\medskip
  \noindent\colorbox{encadre}{%
    \begin{minipage}{\dimexpr\linewidth-2\fboxsep}
      \vspace{2pt}#1\vspace{2pt}
    \end{minipage}}%
  \par\medskip
}
\newcommand{\alerte}[1]{\textcolor{alerte}{\textbf{#1}}}
\newcommand{\bon}[1]{\textcolor{vert}{\textbf{#1}}}

% Chemin de module en chasse fixe : sans point de césure, un nom aussi long que
% media_restorer.engines.triage déborde de 3 cm dans la marge d'un paragraphe
% justifié. \allowbreak après chaque point suffit — le tiret de césure, lui,
% serait trompeur dans un identifiant.
\newcommand{\chemin}[1]{\texttt{#1}}
\newcommand{\modtriage}{\texttt{media\_restorer.\allowbreak engines.\allowbreak triage}}

\newcolumntype{Y}{>{\raggedright\arraybackslash}X}

% =====================================================================
\begin{document}

\begin{center}
  {\LARGE\bfseries\color{accent} Tri grossier du corpus\\[2pt]
   critères globaux robustes et modèles locaux\par}
  \vspace{6pt}
  {\large Rapport préparatoire — projet \texttt{media\_restorer}\par}
  \vspace{4pt}
  {\small Corpus : 11\,313 images, 59,6\,Go \textbullet{} mesures
   \textbf{exhaustives} (11\,310 fichiers, 9\,min\,08\,s de CPU)\par}
  {\small \today\par}
\end{center}

\vspace{4pt}\hrule\vspace{10pt}
\tableofcontents
\vspace{10pt}\hrule

% =====================================================================
\section{Le principe : ce qui rend un critère robuste}

Un critère résiste au changement de domaine (photographie $\rightarrow$ dessin,
puis dessin propre $\rightarrow$ caricature au fusain) quand il repose sur un
signal \textbf{de grande échelle et géométrique}, plutôt que sur une texture ou
une convention de rendu.

\begin{table}[h]
\centering\small
\begin{tabularx}{\textwidth}{YY}
\toprule
\bon{Robuste} — grande échelle, géométrique &
\alerte{Fragile} — fin, dépendant du style \\
\midrule
Nombre de personnages (un / plusieurs / foule) & Identité, nom du personnage \\
Cadrage (gros plan / buste / pied) & Expression, émotion \\
Orientation de la page & Âge, sexe \\
Présence de texte, de bulles & Direction du regard \\
Une case ou plusieurs (bande dessinée) & Couleur des yeux, des cheveux \\
Trait noir ou encre colorée & Type de vêtement \\
Densité d'encre & Pose précise des membres \\
\bottomrule
\end{tabularx}
\caption{La colonne de gauche survit au changement de style parce qu'elle décrit
la \emph{mise en page}, non le \emph{rendu}.}
\end{table}

\aretenir{\textbf{Corollaire opérationnel.} Le tri de premier niveau doit se
limiter à la colonne de gauche. Toute tentative d'y glisser un critère de la
colonne de droite contaminera le tri entier, car rien ne distinguera ensuite un
lot mal classé d'un lot bien classé.}

% =====================================================================
\section{Niveau 0 — les signaux gratuits, sans aucun modèle}
\label{sec:niveau0}

Avant tout modèle, quatre grandeurs se calculent en lisant l'image à résolution
réduite. Les chiffres ci-dessous ne sont pas une estimation : ils proviennent
d'un passage \textbf{exhaustif sur les 11\,310 images} du corpus par le module
\modtriage{} (§\ref{sec:module}).

\begin{table}[h]
\centering\small
\begin{tabular}{llrr}
\toprule
\textbf{Axe} & \textbf{Classe} & \textbf{Effectif} & \textbf{Part} \\
\midrule
Orientation & portrait      & 6\,054 & 53,5\,\% \\
            & paysage       & 4\,594 & 40,6\,\% \\
            & carré         & \phantom{0}\phantom{0}662 & \phantom{0}5,9\,\% \\
\midrule
Couverture d'encre & très clair ($<5\,\%$) & 1\,370 & 12,1\,\% \\
                   & clair (5--15\,\%)     & 3\,944 & 34,9\,\% \\
                   & moyen (15--35\,\%)    & 4\,529 & 40,0\,\% \\
                   & dense ($>35\,\%$)     & 1\,467 & 13,0\,\% \\
\midrule
Résolution & $<1$ Mpx      & 1\,900 & 16,8\,\% \\
           & 1--6 Mpx      & 4\,453 & 39,4\,\% \\
           & 6--20 Mpx     & 3\,288 & 29,1\,\% \\
           & $>20$ Mpx     & 1\,669 & 14,8\,\% \\
\bottomrule
\end{tabular}
\caption{Signaux gratuits, \textbf{corpus entier} (11\,310 images mesurées).
Aucun intervalle de confiance : ce sont des effectifs, pas des estimations.}
\end{table}

\aretenir{\textbf{Sur la taille du corpus.} Le module recense
\textbf{11\,313 fichiers image} là où les décomptes antérieurs en annonçaient
11\,188 : l'écart tient aux 68 \texttt{.webp} et 60 \texttt{.gif} que les
premiers relevés ignoraient. \textbf{3 fichiers} se sont révélés illisibles et
ont été ignorés sans interrompre le parcours — d'où 11\,310 mesures.}

\subsection{Le piège du chromatisme, et sa correction}

Mesurer la saturation \emph{moyenne} d'une image confond deux choses radicalement
différentes : un dessin réellement colorié, et un dessin au trait noir sur un
papier jauni par le temps.

La correction tient en une idée : \textbf{mesurer la saturation séparément sur
les pixels d'encre (les 20\,\% les plus sombres) et sur les pixels de papier
(les 40\,\% les plus clairs).} Le résultat sépare alors deux axes indépendants :

\begin{table}[h]
\centering\small
\begin{tabularx}{\textwidth}{Yrr}
\toprule
\textbf{Combinaison encre / papier} & \textbf{Effectif} & \textbf{Part} \\
\midrule
Trait noir, papier neutre \quad\emph{(dessin au trait « propre »)}
  & 3\,405 & 30,1\,\% \\
\alerte{Trait noir, papier jauni} \quad\emph{(condition de scan, pas contenu)}
  & 3\,807 & 33,7\,\% \\
Encre colorée, papier neutre \quad\emph{(vrai dessin en couleur)}
  & \phantom{0}\phantom{0}535 & \phantom{0}4,7\,\% \\
Encre colorée, papier jauni
  & 3\,563 & 31,5\,\% \\
\bottomrule
\end{tabularx}
\caption{Deux axes valent mieux qu'un : une saturation moyenne unique aurait
rangé les 3\,807 images de la deuxième ligne parmi les dessins « en couleur ».}
\end{table}

\aretenir{\textbf{Deux enseignements pour le corpus.}
\textbf{(1)} \textbf{63,8\,\%} des dessins sont au trait noir (7\,212 images) —
c'est la population homogène sur laquelle un modèle de repères aura le plus de
chances. \textbf{(2)} \textbf{65,2\,\%} ont un papier jauni (7\,370 images) :
c'est un axe de \emph{condition} (candidat naturel à un traitement de
restauration par lot), à ne surtout pas confondre avec un axe de \emph{contenu}.
Ces deux populations se recouvrent largement mais ne coïncident pas — d'où la
nécessité des deux axes.}

\subsection{Coût réel}

Le corpus entier a été traité en \textbf{9 minutes 08 secondes}
(\textbf{48,4\,ms par image}) sur le seul CPU, sans GPU ni téléchargement de
poids.

\alerte{Correction d'une prévision antérieure.} Une première estimation, tirée
d'un échantillon de 500 images, annonçait « 20 à 30 minutes » en redoutant la
queue non-JPEG : \texttt{PIL.draft()} ne réduit le décodage que pour les JPEG,
or le corpus compte \textbf{25,6\,\% de fichiers non-JPEG} (14,5\,\%
\texttt{.png}, 9,9\,\% \texttt{.tif}, plus quelques \texttt{.webp} et
\texttt{.gif}). \textbf{La crainte était infondée} : la mesure exhaustive est
même \emph{plus rapide} par image que l'échantillon (48,4 contre 53\,ms), les
PNG et TIFF du corpus étant en moyenne bien plus petits que les gros scans
JPEG. Le coût reste, dans tous les cas, négligeable.

\subsection{Le module de mesure}
\label{sec:module}

Ces mesures ne sont pas le produit d'un script jetable : elles sont
reproductibles via \modtriage{}, sous-paquet du
c\oe ur sans dépendance Qt (comme \texttt{engines/vectorise} et
\texttt{engines/face\_id}, il ne dérive pas de \texttt{BaseEngine} — le contrat
va d'une image vers des mesures scalaires, non d'une image vers une image).

\begin{itemize}[leftmargin=*,itemsep=3pt]
  \item \texttt{signals.py} — \texttt{ImageSignals} et \texttt{measure\_image()}.
  Les \emph{mesures} sont des attributs, les \emph{catégories} des propriétés
  calculées à la volée depuis des seuils exposés en constantes de module :
  rejouer un tri avec d'autres seuils ne demande pas de relire les fichiers.
  \item \texttt{scan.py} — \texttt{iter\_images()}, \texttt{scan\_directory()},
  \texttt{summarise()}, \texttt{format\_summary()}. La fonction de mesure est
  \textbf{injectable}, ce qui permet de tester parcours, échantillonnage et
  agrégation sans décoder la moindre image ; le rappel \texttt{on\_progress}
  laisse une interface graphique afficher une progression sans que le c\oe ur
  ne connaisse Qt.
\end{itemize}

\alerte{Un piège de reproductibilité mérite d'être signalé :}
\texttt{iter\_images()} \textbf{trie} la liste des fichiers. Sans ce tri,
l'ordre de parcours du système de fichiers rendrait un échantillonnage à graine
fixée non reproductible — vérifié en pratique, deux tirages de 500 images à
graine identique sur liste triée et non triée ne partageaient que
\textbf{21 fichiers sur 500}.

\begin{verbatim}
from media_restorer.engines.triage import (
    scan_directory, summarise, format_summary)

signals = scan_directory("~/Pictures/Dessins")   # ou sample=500 pour calibrer
print(format_summary(summarise(signals)))
\end{verbatim}

% =====================================================================
\section{Niveau 1 — étiqueteurs du domaine illustration}

\subsection{Pourquoi ceux-là et pas un modèle photo}

\texttt{WD14} (\emph{WaifuDiffusion Tagger}) et \texttt{JoyTag} sont entraînés
sur des \textbf{dessins}, selon la taxonomie Danbooru. Le modèle
\texttt{wd-vit-tagger-v3} a été entraîné sur 7,2 millions d'images et ne pèse que
\textbf{94,6 millions de paramètres} — largement à la portée d'un CPU.

Surtout : \textbf{la taxonomie Danbooru possède déjà un vocabulaire canonique
pour exactement les axes recherchés} — nombre de personnages et cadrage sont des
étiquettes de première classe, pas quelque chose à réinventer.

\subsection{Le vocabulaire, trié par transfert attendu}

Le corpus visé (caricatures de presse, fusains, affiches de 1934, numérisations
BNF) n'est pas de l'anime. Il faut donc scinder le vocabulaire selon ce qui
décrit la \emph{géométrie} et ce qui décrit le \emph{style}.

\begin{table}[h]
\centering\small
\begin{tabularx}{\textwidth}{YY}
\toprule
\bon{Devrait transférer} (géométrie, mise en page) &
\alerte{Devrait échouer} (style, contenu) \\
\midrule
\texttt{solo}, \texttt{multiple\_girls}, \texttt{multiple\_boys}, \texttt{6+girls}
  & \texttt{1girl} / \texttt{1boy} : l'attribution de sexe sur une caricature \\
\texttt{full\_body}, \texttt{upper\_body}, \texttt{portrait}, \texttt{close-up}
  & couleur des cheveux, des yeux \\
\texttt{from\_side}, \texttt{profile}
  & vêtements, accessoires \\
\texttt{monochrome}, \texttt{greyscale}
  & noms de personnages \\
\texttt{comic}, \texttt{speech\_bubble}, \texttt{text}
  & \texttt{rating} (classification de contenu) \\
\bottomrule
\end{tabularx}
\caption{Prédiction falsifiable, non simple mise en garde : la colonne de gauche
encode des rapports de surfaces et de positions, indépendants du style de tracé.}
\end{table}

\aretenir{\textbf{La vérification tient en une après-midi :} passer le tagueur sur
100 images, et \textbf{ne noter que la colonne de gauche}. Si elle tient, le tri
de niveau 1 est acquis ; si elle échoue, c'est le niveau 2 qui prend le relais.}

\subsection{Sur le score F1 annoncé}

La fiche du modèle indique \textbf{F1 = 0,4402} au seuil 0,2614. Ce chiffre ne
doit \emph{ni} effrayer \emph{ni} être passé sous silence : c'est une moyenne
multi-étiquettes calculée sur \textbf{des milliers d'étiquettes, en majorité
rares}. Les étiquettes structurelles de tête (\texttt{solo}, \texttt{full\_body}
\dots) — les seules qui nous intéressent — sont bien plus fiables que cette
moyenne ne le laisse croire, précisément parce qu'elles sont fréquentes et
massivement représentées à l'entraînement.

\alerte{Dépendance à ajouter :} \texttt{onnxruntime} (absent), le modèle étant
distribué en ONNX.

% =====================================================================
\section{Niveau 2 — zero-shot CLIP / SigLIP}

\aretenir{\textbf{Ce niveau est essayable dès ce soir, sans installer quoi que ce
soit :} \texttt{transformers} 4.57.6 est déjà une dépendance du projet.
\texttt{onnxruntime}, \texttt{timm}, \texttt{open\_clip} et \texttt{sklearn} sont
tous absents — le niveau 1 demande donc une installation, pas celui-ci.}

Le zero-shot consiste à comparer l'image à des descriptions textuelles et à
retenir la plus proche. Il couvre les axes que Danbooru ne nomme pas, et se
formule directement dans le vocabulaire du corpus :

\begin{itemize}[leftmargin=*,itemsep=3pt]
  \item \emph{« a portrait drawing of a single face »} contre
        \emph{« a full-body drawing of a standing person »} contre
        \emph{« a drawing of a crowd of people »} ;
  \item \emph{« a political cartoon with speech bubbles »} contre
        \emph{« a single figure study »} ;
  \item \emph{« a pencil sketch »} contre \emph{« an ink drawing »} contre
        \emph{« a printed poster »}.
\end{itemize}

Deux propriétés en font le complément naturel du niveau 1 : le vocabulaire est
\textbf{libre} (on écrit les classes que l'on veut, sans réentraîner), et les
mêmes vecteurs d'image serviront ensuite à la sonde linéaire décrite dans le
rapport de faisabilité — l'extraction n'est pas à refaire.

% =====================================================================
\section{Protocole de validation --- la colonne vertébrale}
\label{sec:validation}

C'est le point central, non une annexe. Trier 11\,188 images sur des critères
\emph{réputés} robustes sans les avoir mesurés reviendrait à propager une erreur
inconnue sur tout le corpus.

\begin{enumerate}[leftmargin=*,itemsep=5pt]
  \item \textbf{Échantillon stratifié de 100 images.} Ne pas tirer au hasard
  uniformément : piocher dans chaque sous-dossier thématique
  (\texttt{Cabrol}, \texttt{BNF}, \texttt{Canard}, affiches\dots) pour que les
  familles minoritaires soient représentées.

  \item \textbf{Annoter chaque axe à la main, une seule fois.} Cadrage, nombre de
  personnages, présence de texte, trait noir ou couleur. Une centaine d'images,
  quelques dizaines de minutes.

  \item \textbf{Mesurer l'accord axe par axe} — jamais un score global. Un tri
  peut être excellent sur le cadrage et inutilisable sur le comptage de
  personnages : un score moyen masquerait exactement ce qu'il faut savoir.

  \item \textbf{Ne conserver que les axes au-dessus d'un seuil fixé d'avance.}
  Le choix du seuil appartient à l'utilisateur : 90\,\% pour un tri sur lequel
  s'appuyer aveuglément, 75\,\% pour un tri simplement destiné à guider une revue
  manuelle.

  \item \textbf{Écrire le résultat dans les étiquettes DigiKam.} La plomberie
  existe déjà (\texttt{landmarks.py}, six champs XMP, fusion non destructive) :
  un tri validé devient immédiatement filtrable dans le gestionnaire
  d'étiquettes, et réversible.
\end{enumerate}

\aretenir{\textbf{Ce jeu de 100 images annotées n'est pas jetable.} Il devient le
jeu de validation permanent : tout modèle ultérieur — sonde linéaire, tagueur
affiné, détecteur de repères — se mesure contre lui. C'est le même investissement
qui sert trois fois.}

% =====================================================================
\section{Ordre de mise en \oe uvre}

\begin{table}[h]
\centering\small
\begin{tabularx}{\textwidth}{clYl}
\toprule
& \textbf{Étape} & \textbf{Ce qu'elle apporte} & \textbf{Coût} \\
\midrule
1 & Niveau 0 sur tout le corpus & orientation, encre, papier, résolution & \textbf{9 min CPU} \\
\addlinespace
2 & Échantillon de 100 images annoté & la référence de validation & une soirée \\
\addlinespace
3 & Niveau 2 (zero-shot) & cadrage, comptage, type de planche & aucune dépendance \\
\addlinespace
4 & Niveau 1 (WD14) & les mêmes axes, vocabulaire canonique & \texttt{onnxruntime} \\
\addlinespace
5 & Comparaison 3 \emph{vs} 4 sur l'échantillon & le meilleur des deux, axe par axe & — \\
\bottomrule
\end{tabularx}
\end{table}

L'étape 1 est sans risque et sans dépendance : elle peut être lancée
immédiatement. L'étape 2 conditionne tout le reste. Les étapes 3 et 4 sont
\textbf{interchangeables dans l'ordre} — commencer par celle qui ne demande
aucune installation est simplement plus rapide.

\aretenir{\textbf{Lien avec le rapport de faisabilité.} Ce tri grossier n'est pas
un détour : il produit la population homogène (63,8\,\% de dessins au trait noir,
soit 7\,212 images, dont le sous-ensemble « portrait » ou « buste ») sur
laquelle un détecteur de repères a
le plus de chances de fonctionner. Trier d'abord, c'est aussi choisir les
150 à 300 images à annoter ensuite.}

% =====================================================================
\section*{Sources et références}
\addcontentsline{toc}{section}{Sources et références}

\small\raggedright
\begin{itemize}[leftmargin=*,itemsep=3pt]
  \item \texttt{SmilingWolf/wd-vit-tagger-v3} — fiche du modèle (94,6\,M
  paramètres, 7,2\,M images Danbooru, F1 = 0,4402 au seuil 0,2614, ONNX
  \texttt{onnxruntime} $\geq$ 1.17) —
  \url{https://huggingface.co/SmilingWolf/wd-vit-tagger-v3}
  \item \texttt{JoyTag} — étiqueteur d'illustrations, schéma Danbooru —
  \url{https://github.com/fpgaminer/joytag}
  \item \emph{WaifuDiffusion Tagger} — démonstrateur en ligne —
  \url{https://huggingface.co/spaces/SmilingWolf/wd-tagger}
  \item \texttt{ComfyUI-WD14-Tagger} — référence d'implémentation et
  fonctionnement hors ligne —
  \url{https://github.com/pythongosssss/ComfyUI-WD14-Tagger}
  \item \emph{LP++: A Surprisingly Strong Linear Probe for Few-Shot CLIP},
  CVPR 2024 — pour la suite du niveau 2 —
  \url{https://openaccess.thecvf.com/content/CVPR2024/papers/Huang_LP_A_Surprisingly_Strong_Linear_Probe_for_Few-Shot_CLIP_CVPR_2024_paper.pdf}
\end{itemize}

\medskip
\textbf{Mesures.} Les chiffres de la section \ref{sec:niveau0} proviennent d'un
passage \textbf{exhaustif} du corpus par \modtriage{}
(§\ref{sec:module}), avec \texttt{Pillow} et \texttt{numpy} seuls :
11\,310 images mesurées sur 11\,313 trouvées, en 547,8\,s. Reproductible à
l'identique — le module est couvert par \texttt{tests/test\_triage.py}.

\end{document}
